% ============================================================
% 付録E 市場データおよび市場ショック指標（$\sigma$）の構築方法
% ------------------------------------------------------------
%
% ============================================================

\chapter*{付録 F：市場データおよび市場ショック指標（$\sigma$）の構築方法}
\addcontentsline{toc}{chapter}{付録 F：市場データおよび市場ショック指標（$\sigma$）の構築方法}

本付録では、第8章の実証分析で用いた市場データの取得方法、
市場ショック指標（$\sigma$）の構築手順、
および記述統計の算出方法を詳細に示す。
本研究の透明性と再現性を確保するため、データの出典、加工方法、計算式、
および使用した Python コードを明示する。

\section*{F.1 JEPX スポット価格データの取得方法}
\addcontentsline{toc}{section}{F.1 JEPX スポット価格データの取得方法}

本研究で使用した電力スポット価格データは、
一般社団法人日本卸電力取引所（JEPX）が新サイトにおいて公開している
「スポット市場」データ
（\url{https://www.jepx.jp/electricpower/market-data/spot/}）
から取得した。

JEPX 新サイトでは、以下の情報が日次・30分コマ（1日48コマ）単位で
CSV 形式により提供されている。

\begin{itemize}
  \item 受渡日（YYYY/MM/DD）
  \item 時刻コード（1--48）
  \item システムプライス（円/kWh）
  \item エリアプライス（北海道〜九州）
  \item 売り入札量・買い入札量
  \item 約定総量
\end{itemize}

本研究では、全国共通のシステムプライスを使用した。
データ取得期間は 2016年4月1日から 2025年3月31日までである。

\section*{F.2 日次標準偏差（$\sigma_{\mathrm{daily}}$）の算出方法}
\addcontentsline{toc}{section}{F.2 日次標準偏差（$\sigma_{\mathrm{daily}}$）の算出方法}

1日48コマのシステムプライスを用いて、日次標準偏差を以下の式で算出した。



\[
\sigma_{\mathrm{daily}}(t)
=
\mathrm{StdDev}
\left(
\mathrm{SystemPrice}_{t,1},\ldots,\mathrm{SystemPrice}_{t,48}
\right)
\]



異常値（マイナス価格など）が存在する場合は、その日のデータを除外した。

\section*{F.3 月次市場ショック（$\sigma_{\mathrm{monthly}}$）の算出方法}
\addcontentsline{toc}{section}{F.3 月次市場ショック（$\sigma_{\mathrm{monthly}}$）の算出方法}

市場ショック指標は、日次標準偏差の月次標準偏差として定義した。



\[
\sigma_{\mathrm{monthly}}
=
\mathrm{StdDev}\left(\sigma_{\mathrm{daily}} \in \text{month}\right)
\]



この指標は、当該月における価格変動の激しさを表すものであり、
新電力の調達コストの不確実性を捉える proxy として妥当である。

\section*{F.4 価格データの例示}
\addcontentsline{toc}{section}{F.4 価格データの例示}

表 E-1 に、2021年1月と2020年平常月の
システムプライスの例を示す。

\begin{table}[H]
\centering
\caption{システムプライスの例（抜粋）}
\begin{tabular}{lcc}
\toprule
日付 & 平均価格（円/kWh） & 日次標準偏差 \\
\midrule
2021/01/08 & 154.2 & 87.5 \\
2021/01/12 & 197.3 & 102.4 \\
2020/10/15 & 7.8 & 2.1 \\
2020/10/20 & 8.3 & 1.9 \\
\bottomrule
\end{tabular}
\end{table}

2021年1月の価格変動が極端に大きいことが確認できる。

\section*{F.5 市場ショック指標（$\sigma$）の推移}
\addcontentsline{toc}{section}{F.5 市場ショック指標（$\sigma$）の推移}

図 E-1 に、2016年4月〜2025年3月の
市場ショック指標（$\sigma$）の推移を示す。

\begin{figure}[H]
\centering
\includegraphics[width=14cm]{D:/texlive/2026/work/thesis_Rev1_20260516/figures/sigma_trend.pdf}
\caption{市場ショック指標（$\sigma$）の推移}
\end{figure}

2021年1月に突出したピークが観察される。

\section*{F.6 データ利用上の留意点}
\addcontentsline{toc}{section}{F.6 データ利用上の留意点}

\begin{itemize}
  \item JEPX 新サイトは旧サイトと異なり、1年分のデータを一括取得できるため、欠損が少ない。
  \item 本研究では、全国共通のシステムプライスを採用した。
        エリアプライスは地域差が大きいが、倒産企業の多くが複数エリアで調達していたためである。
  \item データ加工は日次標準偏差および月次集計に限定し、恣意的な補正は行っていない。
\end{itemize}

\section*{F.7 記述統計の算出方法}
\addcontentsline{toc}{section}{F.7 記述統計の算出方法}

第8章の記述統計（表 8-1）は、
panel\_data\_with\_sigma.csv に含まれる
4 変数（$\sigma$, $S$, $D$, $\text{default}$）を対象とし、
Python（pandas）により算出した。

使用したコードは以下の通りである。

\begin{verbatim}
import pandas as pd
import os

DATA_DIR = r"D:\texlive\2026\work\figure\JEPX_Spot_Price_Data"
panel_path = os.path.join(DATA_DIR, "panel_data_with_sigma.csv")

df = pd.read_csv(panel_path, encoding="cp932")

vars = ["sigma", "S", "D", "default"]
desc = df[vars].describe()
print(desc)
\end{verbatim}

算出した統計量は、平均値・標準偏差・最小値・最大値であり、
欠損値は dropna() により除去した。
これらの値を第8章の記述統計表に反映した。

